% ===========================================================================
%  樊启斌《高等代数典型问题与方法》 · 第 2 章 行列式 —— 深化提取（服务于
%  「余子式和代数余子式的计算」，对应考研大纲「行列式」）
%
%  源：PDF p114–p121 = 印刷页 105–112（实测顶部页码：p114=105、p115=106、
%      p116=107、p117=108、p118=109、p119=110、p120=111、p121=112，
%      即偏移 = PDF 页 = 印刷页 + 9，与 AGENTS.md §3.4 一致）
%
%  覆盖：§2.3 代数余子式问题（印刷页 105–108，全，★本次核心目标）
%        §2.4 综合性问题（印刷页 109–112，默认整节跳过；仅收 1 条考纲内
%        可用的通用结论「对全排列求和的抽象行列式必为 0」，见例 2.64 一块）
%
%  提取原则：只取定义 / 定理 / 性质 / 重要结论 / 方法要点 / 关键证明思路。
%            例题的数字演算过程一律不进正文，只在承载了可复用结论时提取
%            「结论 + 一句话方法」。
% ===========================================================================


% ---------------------------------------------------------------------------
% 【深化】服务考点：代数余子式的定义与基本展开式（考纲「行列式」）
% 来源：樊 2.3，印刷页 105（PDF p114）【解】段
% 反哺点：把"余子式之和"问题统一落到一个升阶行列式上，一行一列地看就行
% ---------------------------------------------------------------------------
\begin{fdeep}{余子式的两个基本计算式（升阶法与降阶法）}
\fsec{§2.3\quad 代数余子式问题}

设 $D$ 是 $n$ 阶行列式，$M_{ij}$ 是元素 $a_{ij}$ 的余子式，则

（1）\textbf{降阶式}（把某一行换成"和为 $1$"的行）：
\fml{\begin{vmatrix}
\boxed{1} & 1 & \cdots & 1\\
a_{21} & a_{22} & \cdots & a_{2n}\\
\vdots & \vdots & & \vdots\\
a_{n1} & a_{n2} & \cdots & a_{nn}
\end{vmatrix}=M_{11}+M_{12}+\cdots+M_{1n}.}

（2）\textbf{升阶式}（在 $D$ 的最上方加一行 $1\,0\,\cdots\,0$、最左方加一列，得 $D$ 的升阶行列式）：
\fml{\begin{vmatrix}
1 & 0 & 0 & \cdots & 0\\
0 & 1 & 1 & \cdots & 1\\
1 & a_{21} & a_{22} & \cdots & a_{2n}\\
\vdots & \vdots & \vdots & & \vdots\\
1 & a_{n1} & a_{n2} & \cdots & a_{nn}
\end{vmatrix}=\begin{vmatrix}
1 & a_{11} & \cdots & a_{1n}\\
1 & a_{21} & \cdots & a_{2n}\\
\vdots & \vdots & & \vdots\\
1 & a_{n1} & \cdots & a_{nn}
\end{vmatrix}.}

升阶后把第 $1$ 列的 $a_{ij}$ 倍加到第 $j+1$ 列（$j=1,2,\cdots,n$）即得上式右端；再按第一列展开，便得到
\fml{D=\begin{vmatrix}
1 & 1 & \cdots & 1\\
a_{21} & a_{22} & \cdots & a_{2n}\\
\vdots & \vdots & & \vdots\\
a_{n1} & a_{n2} & \cdots & a_{nn}
\end{vmatrix}+\begin{vmatrix}
a_{11} & a_{12} & \cdots & a_{1n}\\
1 & 1 & \cdots & 1\\
\vdots & \vdots & & \vdots\\
a_{n1} & a_{n2} & \cdots & a_{nn}
\end{vmatrix}+\cdots+(-1)^{n}\begin{vmatrix}
a_{11} & a_{12} & \cdots & a_{1n}\\
\vdots & \vdots & & \vdots\\
a_{n-1,1} & a_{n-1,2} & \cdots & a_{n-1,n}\\
1 & 1 & \cdots & 1
\end{vmatrix}.}
其中第一个行列式按第一行展开，其余各个行列式都按第 $1$ 行展开即为 $M$ 之和。

\textbf{【反哺点】}凡是要算"某一行（列）的余子式之和"，不必逐个求 $n$ 个 $n-1$ 阶行列式：直接把那一行（列）换成全 $1$，或把 $D$ 的一行（列）拆成"$D$ 的对应行 $+$ 全 $1$ 行"，一个 $n$ 阶行列式即得答案。
\end{fdeep}


% ---------------------------------------------------------------------------
% 【深化】服务考点：代数余子式的性质——把所有元素加上同一个数（考纲「行列式」）
% 来源：樊 2.3 例 2.55（浙大 2006 / 西北大学 2015），印刷页 105–106（PDF p114–115）
% 反哺点：$\sum\sum A_{ij}$ 类问题有了闭式公式，$|A+xJ|$ 也可一步写出
% ---------------------------------------------------------------------------
\begin{fdeep}{所有元素都加上同一个数 $x$ 后的行列式与 $\sum\sum A_{ij}$}
\textbf{结论（例 2.55）}\quad 记 $A=(a_{ij})_{n\times n}$，$A_{ij}$ 是 $a_{ij}$ 的代数余子式，$J$ 为元素全为 $1$ 的 $n$ 阶方阵，则
\fml{\begin{vmatrix}
a_{11}+x & a_{12}+x & \cdots & a_{1n}+x\\
a_{21}+x & a_{22}+x & \cdots & a_{2n}+x\\
\vdots & \vdots & & \vdots\\
a_{n1}+x & a_{n2}+x & \cdots & a_{nn}+x
\end{vmatrix}=|A|+x\sum_{i=1}^{n}\sum_{j=1}^{n}A_{ij},}
即 $|A+xJ|=|A|+x\sum\limits_{i=1}^{n}\sum\limits_{j=1}^{n}A_{ij}$。特别地，$x=1$ 时
\fml{|A+J|=|A|+\sum_{i=1}^{n}\sum_{j=1}^{n}A_{ij}.}

\textbf{推导要点}\quad 只需三个动作，全程只用行列式的线性性质：
\fml{\begin{aligned}
|A+xJ|&=\begin{vmatrix}
a_{11}+x & a_{12}+x & \cdots & a_{1n}+x\\
a_{21} & a_{22} & \cdots & a_{2n}\\
\vdots & \vdots & & \vdots\\
a_{n1} & a_{n2} & \cdots & a_{nn}
\end{vmatrix}+\begin{vmatrix}
x & x & \cdots & x\\
a_{21} & a_{22} & \cdots & a_{2n}\\
\vdots & \vdots & & \vdots\\
a_{n1} & a_{n2} & \cdots & a_{nn}
\end{vmatrix}\\
&=\ \cdots\ =\ |A|+x\begin{vmatrix}
1 & 1 & \cdots & 1\\
a_{21} & a_{22} & \cdots & a_{2n}\\
\vdots & \vdots & & \vdots\\
a_{n1} & a_{n2} & \cdots & a_{nn}
\end{vmatrix}
\end{aligned}}
（对每一行都用"把第一行乘 $-1$ 加到其余各行，再把第一行拆项"，逐行做 $n$ 次），
最后那个行列式恰好是 $\sum\limits_{i=1}^{n}\sum\limits_{j=1}^{n}A_{ij}$。

\textbf{顺带得到的两件事}

（a）\textbf{余子式之和与 $x$ 无关}\quad 把 $D=|A+xJ|$ 中 $(i,j)$ 元的代数余子式记为 $D_{ij}$，则
\fml{\sum_{i=1}^{n}\sum_{j=1}^{n}D_{ij}=\sum_{i=1}^{n}\sum_{j=1}^{n}A_{ij}}
（因为 $|A|=D-x\sum\sum D_{ij}$，代回上式即得），故 $D$ 的所有元素的代数余子式之和与 $x$ 无关。

（b）\textbf{把 $x$ 取成 $-a_{11}$ 等具体值}，可把"含参行列式"化为"求 $\sum\sum A_{ij}$"；
反过来，$\sum\sum A_{ij}$ 又等于
\fml{\sum_{i=1}^{n}\sum_{j=1}^{n}A_{ij}=|A|\cdot\sum_{i=1}^{n}\sum_{j=1}^{n}\bigl(A^{-1}\bigr)_{ij}\qquad (|A|\neq0),}
即 $\sum\sum A_{ij}$ 与 $A^{-1}$ 的全部元素之和只差一个因子 $|A|$。

\textbf{【反哺点】}考纲内一旦出现"$\sum\sum A_{ij}$"或"所有元素同加一个数"，直接套 $|A+xJ|=|A|+x\sum\sum A_{ij}$：把已知行列式凑成 $A+xJ$ 的形状，未知的 $\sum\sum A_{ij}$ 就是一次多项式的一次项系数。
\end{fdeep}


% ---------------------------------------------------------------------------
% 【深化】服务考点：$\sum A_{i1}$、$\sum A_{ij}$ 的通用求法（考纲「行列式」）
% 来源：樊 2.3 例 2.55(2) 的证明原理，印刷页 105–106（PDF p114–115）
% 反哺点：三种"余子式之和"用同一个动作解决——加一行（列）全 1
% ---------------------------------------------------------------------------
\begin{fdeep}{余子式之和、代数余子式之和的通用求法（加全 $1$ 行/列）}
\textbf{统一原理}\quad $A_{ij}$ 与 $a_{ij}$ 的取值无关。于是：
\begin{itemize}
\item 在 $A$ 的下方添一行全 $1$，按这一行展开，得
\fml{\sum_{j=1}^{n}A_{ij}=\begin{vmatrix}
a_{11} & a_{12} & \cdots & a_{1n}\\
\vdots & \vdots & & \vdots\\
a_{i-1,1} & a_{i-1,2} & \cdots & a_{i-1,n}\\
1 & 1 & \cdots & 1\\
a_{i+1,1} & a_{i+1,2} & \cdots & a_{i+1,n}\\
\vdots & \vdots & & \vdots\\
a_{n1} & a_{n2} & \cdots & a_{nn}
\end{vmatrix};}
\item 在 $A$ 的右方添一列全 $1$，按这一列展开，得 $\sum\limits_{i=1}^{n}A_{ij}$；
\item 添一行全 $1$ \emph{且}添一列全 $1$（即 $\begin{vmatrix} A & \mathbf{1}\\ \mathbf{1}^{T} & 1\end{vmatrix}$），按最后一行展开，得 $\sum\limits_{i=1}^{n}\sum\limits_{j=1}^{n}A_{ij}$。
\end{itemize}

\textbf{三种读法（同一件事）}\quad 对元素全为 $1$ 的 $n$ 阶方阵 $J$：
\fml{\sum_{i=1}^{n}\sum_{j=1}^{n}A_{ij}=\bigl|A+J\bigr|-|A|=\bigl|A\bigr|\cdot\mathbf{1}^{T}A^{-1}\mathbf{1}
=\Bigl|A\bigr|\sum_{i=1}^{n}\sum_{j=1}^{n}\bigl(A^{-1}\bigr)_{ij}.}
其中中间一式来自 $A^{*}=|A|A^{-1}$ 与 $\mathbf{1}^{T}A^{*}\mathbf{1}=\sum\sum A_{ij}$——即"$\sum\sum A_{ij}$ 就是用全 $1$ 向量对伴随矩阵做一次双线性型"。

\textbf{【反哺点】}考纲内三种问法全归一个动作：看到 $\sum\limits_{i=1}^{n}A_{i1}$、$\sum\limits_{j=1}^{n}A_{1j}$、$\sum\sum A_{ij}$，就在相应位置补全 $1$ 的行/列，化成 $n$ 阶（或 $n+1$ 阶）行列式；反之给出 $\sum\sum A_{ij}$ 也能反解 $|A+J|$。
\end{fdeep}


% ---------------------------------------------------------------------------
% 【深化】服务考点：代数余子式的线性组合（考纲「行列式」+「矩阵」伴随矩阵）
% 来源：樊 2.3 例 2.59，印刷页 108（PDF p117）
% 反哺点：把"换一行的行列式之和"降成 $\sum A_{ij}$ 的线性组合，是考纲内 k 倍行加法的逆用法
% ---------------------------------------------------------------------------
\begin{fdeep}{代数余子式的线性组合：行替换公式与列替换公式}
\textbf{行替换公式}\quad 设 $n$ 阶行列式
\fml{D=\begin{vmatrix}
a_{11} & a_{12} & \cdots & a_{1,n-1} & 1\\
a_{21} & a_{22} & \cdots & a_{2,n-1} & 1\\
\vdots & \vdots & & \vdots & \vdots\\
a_{n1} & a_{n2} & \cdots & a_{n,n-1} & 1
\end{vmatrix},}
把 $D$ 的第 $i$ 行换成 $x_1,x_2,\cdots,x_{n-1},1$、其余行不动，所得行列式记为 $D_i\ (i=1,2,\cdots,n)$，则
\fml{D_1+D_2+\cdots+D_n=D+x_1\sum_{i=1}^{n}A_{i1}+x_2\sum_{i=1}^{n}A_{i2}+\cdots+x_{n-1}\sum_{i=1}^{n}A_{i,n-1}.}
\textbf{证明就是"按行展开"}\quad $D_i$ 按它被换掉的那一行展开：
\fml{D_i=x_1A_{i1}+x_2A_{i2}+\cdots+x_{n-1}A_{i,n-1}+1\cdot A_{in}\qquad(i=1,2,\cdots,n),}
对 $i$ 求和，注意 $\sum\limits_{i=1}^{n}A_{in}=D$，即得结论。

\textbf{两条派生的常用特例}
\begin{itemize}
\item 若取 $x_1=x_2=\cdots=x_{n-1}=0$（即把第 $i$ 行换成 $0,\cdots,0,1$），则 $D_i=A_{in}$，故 $\sum\limits_{i=1}^{n}A_{in}=D$——这正是"按最后一列展开"。
\item 若取 $x_1=1,x_2=\cdots=x_{n-1}=0$（把第 $i$ 行换成 $1,0,\cdots,0,1$），则因第 $1$ 列与最后一列成比例，行列式为 $0$，于是
\fml{\sum_{i=1}^{n}A_{i1}=\begin{vmatrix}
1 & a_{12} & \cdots & a_{1,n-1} & 1\\
1 & a_{22} & \cdots & a_{2,n-1} & 1\\
\vdots & \vdots & & \vdots & \vdots\\
1 & a_{n2} & \cdots & a_{n,n-1} & 1
\end{vmatrix}=0.}
\end{itemize}

\textbf{列替换公式（对偶形式，同样成立）}\quad 行列式等于任一行的元素与其代数余子式乘积之和，因此把第 $j$ 列换成 $y_1,y_2,\cdots,y_n$，所得行列式的值为
\fml{y_1A_{1j}+y_2A_{2j}+\cdots+y_nA_{nj}\qquad(j=1,2,\cdots,n).}

\textbf{【反哺点】}考纲内"某一行换掉后得到的行列式之和"这类题，不要一个个算：$D_i$ 按被换的那一行展开，立刻化为 $\sum\limits_{i=1}^{n}A_{ij}$ 的线性组合；反过来，知道 $\sum\sum A_{ij}$ 也能读出"换行行列式"的值。
\end{fdeep}


% ---------------------------------------------------------------------------
% 【深化】服务考点：代数余子式的性质（考纲「行列式」+「矩阵」伴随矩阵）
% 来源：樊 2.3 例 2.57，印刷页 107（PDF p116）
% 反哺点：给出一类"所有代数余子式全相等"的判别条件，直接服务于 $\sum\sum A_{ij}$ 的求值
% ---------------------------------------------------------------------------
\begin{fdeep}{行和、列和全为 $0$ 时所有代数余子式相等}
\textbf{结论（例 2.57）}\quad 设 $A=(a_{ij})_{n\times n}$ 满足
\fml{\begin{cases}
a_{i1}+a_{i2}+\cdots+a_{in}=0,\\[1mm]
a_{1j}+a_{2j}+\cdots+a_{nj}=0
\end{cases}\qquad(i,j=1,2,\cdots,n),}
（即 $A$ 的每一行、每一列元素之和都为 $0$）则 $A$ 的所有元素的代数余子式相等，即
\fml{A_{ij}=A_{11}\qquad(i,j=1,2,\cdots,n).}

\textbf{关键证明思路（$A_{ij}=A_{i1}$ 的一段）}\quad 对第一行的代数余子式，$j=2,3,\cdots,n$ 时
\fml{A_{1j}=(-1)^{1+j}\begin{vmatrix}
a_{21} & \cdots & a_{2,j-1} & a_{2,j+1} & \cdots & a_{2n}\\
a_{31} & \cdots & a_{3,j-1} & a_{3,j+1} & \cdots & a_{3n}\\
\vdots & & \vdots & \vdots & & \vdots\\
a_{n1} & \cdots & a_{n,j-1} & a_{n,j+1} & \cdots & a_{nn}
\end{vmatrix}.}
把第 $2$ 列至第 $n-1$ 列都加到第 $1$ 列，利用 $\sum\limits_{j=1}^{n}a_{ij}=0\ (i=2,3,\cdots,n)$ 得
\fml{A_{1j}=(-1)^{1+j}\begin{vmatrix}
-a_{2j} & a_{22} & \cdots & a_{2,j-1} & a_{2,j+1} & \cdots & a_{2n}\\
-a_{3j} & a_{32} & \cdots & a_{3,j-1} & a_{3,j+1} & \cdots & a_{3n}\\
\vdots & \vdots & & \vdots & \vdots & & \vdots\\
-a_{nj} & a_{n2} & \cdots & a_{n,j-1} & a_{n,j+1} & \cdots & a_{nn}
\end{vmatrix},}
再把第 $1$ 列的负号提出，并把第 $1$ 列依次与第 $2$ 列、第 $3$ 列、$\cdots$、第 $j-1$ 列互换（共 $j-2$ 次），得
\fml{A_{1j}=(-1)^{1+j}(-1)^{1}\cdot(-1)^{j-2}\begin{vmatrix}
a_{22} & \cdots & a_{2,j-1} & a_{2j} & a_{2,j+1} & \cdots & a_{2n}\\
a_{32} & \cdots & a_{3,j-1} & a_{3j} & a_{3,j+1} & \cdots & a_{3n}\\
\vdots & & \vdots & \vdots & \vdots & & \vdots\\
a_{n2} & \cdots & a_{n,j-1} & a_{nj} & a_{n,j+1} & \cdots & a_{nn}
\end{vmatrix}=A_{11}.}
对第 $i\ (i=2,3,\cdots,n)$ 行的代数余子式 $A_{ij}$ 重复上述过程，得 $A_{ij}=A_{i1}$；同理用 $a_{1j}+a_{2j}+\cdots+a_{nj}=0$ 可证 $A_{i1}=A_{11}$。故 $A_{ij}=A_{11}$。

\textbf{【反哺点】}遇到"行和与列和都为 $0$"的矩阵，$\sum\sum A_{ij}=n^{2}A_{11}$，只需算一个代数余子式；这也是判断"某矩阵的伴随矩阵是否为常矩阵（数量矩阵）"的最快入口。
\end{fdeep}


% ---------------------------------------------------------------------------
% 【深化】服务考点：行列式的性质（转置不变、乘法）与代数余子式（考纲「行列式」「矩阵」）
% 来源：樊 2.3 例 2.60，印刷页 109（PDF p118）
% 反哺点：$A_{ij}=a_{ij}$ 的实矩阵必可逆，且给出 $|A|^{n-2}=1$，直接服务伴随矩阵求逆
% ---------------------------------------------------------------------------
\begin{fdeep}{$A_{ij}=a_{ij}$ 的实矩阵：可逆性与 $|A|^{n-2}=1$}
\textbf{结论（例 2.60，上海师范大学 2002）}\quad 设 $a_{ij}\in\mathbb{R}$（实矩阵 $A=(a_{ij})_{n\times n}$），且 $A$ 中至少有一个元素 $a_{ij}\neq0$。若 $A$ 的一切元素 $a_{ij}$ 的代数余子式 $A_{ij}=a_{ij}$，则
\fml{D^{n-2}=1\qquad\bigl(D=|A|\bigr).}

\textbf{思路}\quad 取 $a_{ij}\neq0$，按其所在列展开：
\fml{D=\sum_{i=1}^{n}a_{ij}A_{ij}=\sum_{i=1}^{n}a_{ij}^{2}>0;}
又因 $A_{ij}=a_{ij}$ 即 $A^{*}=A^{T}$，由行列式乘法规则
\fml{D^{2}=DD^{T}=\begin{vmatrix}
a_{11} & a_{12} & \cdots & a_{1n}\\
a_{21} & a_{22} & \cdots & a_{2n}\\
\vdots & \vdots & & \vdots\\
a_{n1} & a_{n2} & \cdots & a_{nn}
\end{vmatrix}\begin{vmatrix}
A_{11} & A_{21} & \cdots & A_{n1}\\
A_{12} & A_{22} & \cdots & A_{n2}\\
\vdots & \vdots & & \vdots\\
A_{1n} & A_{2n} & \cdots & A_{nn}
\end{vmatrix}=\begin{vmatrix}
D & 0 & \cdots & 0\\
0 & D & \cdots & 0\\
\vdots & \vdots & & \vdots\\
0 & 0 & \cdots & D
\end{vmatrix}=D^{n},}
故 $D^{n-2}=1$。

\textbf{【反哺点】}考纲内凡是给"代数余子式与元素本身的关系"，都翻译成伴随矩阵语言 $A^{*}=A^{T}$（或 $A^{*}=f(A)$），再用 $AA^{*}=|A|E$ 一次算出 $|A|$；本题还顺带说明：实矩阵不必对称，只要 $A_{ij}=a_{ij}$ 就一定有 $|A|>0$。
\end{fdeep}


% ---------------------------------------------------------------------------
% 【深化】服务考点：行列式的性质（交换两列变号 / 展开定理）（考纲「行列式」）
% 来源：樊 2.4 例 2.64（中山大学 2008），印刷页 111（PDF p120）
% 说明：§2.4 默认整节跳过；此条是"对全排列求和"的通用结论，且只用到考纲内的
%       "交换两列行列式变号"，故单独提取。
% 反哺点：含 $\sum$ 与排列的抽象行列式，用一次列交换即可判断其值为 0
% ---------------------------------------------------------------------------
\begin{fdeep}{对全排列求和的抽象行列式必为 $0$}
\textbf{结论（例 2.64，中山大学 2008）}\quad 设 $A=(a_{ij})_{n\times n}$，则
\fml{\sum_{j_1j_2\cdots j_n}\begin{vmatrix}
a_{1j_1} & a_{1j_2} & \cdots & a_{1j_n}\\
a_{2j_1} & a_{2j_2} & \cdots & a_{2j_n}\\
\vdots & \vdots & & \vdots\\
a_{nj_1} & a_{nj_2} & \cdots & a_{nj_n}
\end{vmatrix}=0,}
这里 $\sum$ 是对 $1,2,\cdots,n$ 的所有全排列 $j_1j_2\cdots j_n$ 求和。

\textbf{思路}\quad 记 $\tau(j_1j_2\cdots j_n)$ 表示排列 $j_1j_2\cdots j_n$ 的逆序数，由行列式交换两列变号知
\fml{\begin{vmatrix}
a_{1j_1} & a_{1j_2} & \cdots & a_{1j_n}\\
a_{2j_1} & a_{2j_2} & \cdots & a_{2j_n}\\
\vdots & \vdots & & \vdots\\
a_{nj_1} & a_{nj_2} & \cdots & a_{nj_n}
\end{vmatrix}=(-1)^{\tau(j_1j_2\cdots j_n)}\det A,}
而奇排列与偶排列的个数相等，都等于 $\dfrac{n!}{2}$，故求和为 $0$。

\textbf{【反哺点】}考纲内遇到"对全部 $n!$ 个排列求和"的抽象行列式，不要展开：把每个排列对应的行列式化为 $(-1)^{\tau}\det A$，利用奇偶排列成对相消；这一步只用到"交换两列变号"，属于行列式性质的标准应用。
\end{fdeep}


% ---------------------------------------------------------------------------
% 以下按深化提取原则跳过
% ---------------------------------------------------------------------------
% fan p106 例 2.56（$n^{2}$ 个互异数的适当排列使 $n$ 阶行列式非零）：归纳法存在性
%          证明，结论与考纲"行列式的计算"无接口，仅作理解，不提取。
% fan p107–108 例 2.58（北京工业大学 2012，$f(X)=X^{T}A^{*}X$ 与 $|A_{3}|$）：
%          其可复用内核"$\sum\sum A_{ij}=|A+J|-|A|$"已并入上文
%          「所有元素都加上同一个数 $x$」一条，故不另立块。
% fan p109 例 2.61（最大公因数 $r$ 整除行列式 $d$）：数论背景，与考纲无关，跳过。
% fan p110 例 2.62（$k$ 行与 $j$ 列交叉处元素全为 $0$ 且 $k+j>n$ $\Rightarrow D=0$）：
%          §2.4 竞赛级结论，考纲内无对应题型，跳过。
% fan p110–111 例 2.63（华中师范大学 1994，$D_{n}$ 展开式中正项项数 $x=\frac12(2^{n-1}+n!)$）：
%          虽用到"正负项个数之差 $=$ 行列式值"这一技巧，但结论针对特定行列式，跳过。
% fan p111 例 2.65（北京大学 2018，$3$ 阶 $(0,1)$ 行列式最大值 $2$）：
%          竞赛级最值问题，超出考纲，跳过。
% fan p112 例 2.66（华东师范大学 2005，元素全为 $\pm1$ 时 $|D_{n}|\leqslant(n-1)(n-1)!$）
%          及其【注】$|D_{n}|\leqslant\frac{2}{3}n!$：竞赛级不等式，超出考纲，跳过。
% fan p112 例 2.67（中国科学技术大学 2007，对角元为正、非对角元为负、列和为正
%          $\Rightarrow\det A>0$）：竞赛级，超出考纲，跳过。
% ===========================================================================
